\documentclass[10pt,a4paper]{amsart}
\usepackage[margin=19mm]{geometry}
\usepackage{amsmath,amssymb,mathtools}
\usepackage[scr=rsfs,cal=euler]{mathalfa}
\usepackage{xfrac}
\usepackage[unicode,hidelinks,bookmarks=false]{hyperref}
\newcommand{\ie}{i.e.\ }
\newcommand{\cf}{cf.\ }
\newcommand{\resp}{respectively\ }
\newcommand{\Subsection}{\subsection}
\providecommand{\nobreakdash}{\nobreak\hbox{-}}
\setlength{\emergencystretch}{2em}
\multlinegap0pt
\usepackage{bm}
\usepackage{bbm}
\usepackage{dsfont}
\usepackage{xspace}
\usepackage{cancel}
\usepackage{mathtools}
\usepackage{amsthm}
\makeatletter
\@ifundefined{theorem}{\newtheorem{theorem}{Theorem}}{}
\@ifundefined{lemma}{\newtheorem{lemma}[theorem]{Lemma}}{}
\makeatother
\title[Precise propagation profile for some monostable free boundar]{Precise propagation profile for some monostable free boundary problems in time-periodic media}
\author{Yihong Du, Zhuo Ma, Zhi-Cheng Wang}
\date{Source: arXiv:2602.05328v1 (CC BY 4.0)}
\begin{document}
\maketitle

\noindent\emph{Source and licence.} Precise propagation profile for some monostable free boundary problems in time-periodic media. Version arXiv:2602.05328v1 (CC BY 4.0), distributed under the Creative Commons Attribution 4.0 International licence, as recorded in the arXiv licence metadata for this exact version and shown on the abstract page https://arxiv.org/abs/2602.05328v1. Full rights notice: licenses/arxiv-2602.05328-CC-BY-4.0.txt.\par\smallskip

\noindent\textit{Editorial scope.} Counted unit: the Lemma whose source label is \texttt{lemm-gr} in the article (source0.tex lines 1298--1300, source label \texttt{lemm-gr}), delivered together with the article's own complete original proof of it (source0.tex lines 1301--1516, one \texttt{proof} environment of 216 lines). No mathematics is written, repaired or re-derived: statement and proof are the article's own text line for line. Required context retained uncounted: lem-lim (lines 1175--1188); eq-vpr (lines 1292--1296). Every definition, display and label of those lines is retained unchanged. Omitted from this package and not redistributed: 1--1174, 1189--1291, 1297--1297, 1517--1850. Nothing inside a retained excerpt is omitted or altered: every retained body is delivered exactly as the article's lines are written, so no omission receipt of any kind accompanies this package. Citations: the retained excerpts carry no citation command at all, so no bibliography entry of the article is transcribed and no reference is redistributed. Reference adaptation: none was needed and none was made; every reference inside the delivered unit resolves inside it, so no unresolved reference or citation is produced. Printed slips kept literal: no printed word, symbol, hypothesis or number of the counted statement or of its proof was altered. Layout and class: the article's own class is \texttt{amsart} with packages including amsfonts, bm, mathrsfs, hyperref, color, cite, enumerate; the wrapper is the lane's pinned \texttt{amsart} editorial wrapper at 10pt on A4, which restates the theorem-like environments the retained text needs and adds the article's own macro definitions that the retained text needs (none), with extra wrapper packages (bm, bbm, dsfont, xspace, cancel, mathtools). The delivered numbering is the unit's own. Assets: no retained span contains a figure environment, a graphics-inclusion command, a PDF-page inclusion, a table or a TikZ picture; no image file is copied into this package, the package declares no asset dependency and no image file is redistributed with it. Licence: version arXiv:2602.05328v1 of this article is distributed under the Creative Commons Attribution 4.0 International licence, as recorded in the arXiv licence metadata for this exact version and shown on the abstract page https://arxiv.org/abs/2602.05328v1. A full rights notice is delivered at licenses/arxiv-2602.05328-CC-BY-4.0.txt. Characters: every retained excerpt is pure ASCII, so the delivered engine, XeLaTeX with its default Latin Modern text font, reports no missing character.
\begin{lemma}\label{lem-lim}
  There exists a subsequence of $\{(\rho_n, v_n)\}$, still denoted by  itself, such that
  \begin{equation*}
    \rho_n\to \Gamma \text{ in } C^{1+\alpha/2}_{loc}(\mathbb{R}) \text{  and }\|v_n-V\|_{C_{loc}^{(1+\alpha)/2,1+\alpha}(\overline\Sigma)}\to 0 ~\text{ as } n\to \infty,
  \end{equation*}
  where $\Sigma:=\{(t,x):t\in\mathbb{R}, -\infty<x< \Gamma(t)\}$. Moreover $(V(t,x),\Gamma(t))$ satisfies
  \begin{equation}\label{eq-V}
    \begin{cases}
      V_t-dV_{xx}-k^*(t)V_x=f(t,V), ~V>\delta, &(t,x)\in\Sigma,\\
      V(t,\Gamma(t))=\delta, &t\in\mathbb{R},\\
      \Gamma'(t)=-\frac{d}{\delta}V_x(t,\Gamma(t))-k^*(t),~ \Gamma(t)\geq \Gamma(r_*), &t\in\mathbb R.
    \end{cases}
  \end{equation}
\end{lemma}

\begin{equation}\label{eq-vpr}
 \begin{cases} V(t,x)\geq \Phi^*(t,L^*-x) \text{ for }(t,x)\in\mathbb{R}\times(-\infty,L^*],\\
 \min_{t\in\mathbb R}\Gamma(t)=\Gamma (r_*)\geq  L^*.
 \end{cases}
\end{equation}

  \begin{lemma}\label{lemm-gr}
$L^*=\Gamma(r_*)$.
\end{lemma}

\begin{proof}
Suppose to the contrary that $L^*<\Gamma(r_*)$. We are going to derive a contradiction in  several steps.

\textbf{ Step 1.} We prove that 
\begin{equation}\label{eq-v>p}
  V(t,x)> \Phi^*(t,L^*-x) ~~\text{ for } (t,x)\in\mathbb R \times(-\infty,L^*].
\end{equation}
Otherwise, in view of \eqref{eq-vpr}, there exists $(t_0,x_0)\in \mathbb R \times(-\infty,L^*]$ such that
\begin{equation*}
  V(t_0,x_0)= \Phi^*(t_0,L^*-x_0).
\end{equation*}
Since $L^*<\Gamma(r_*)\leq \Gamma(t)$, we obtain $V(t,L^*)>\delta=\Phi^*(t,0)$ for all $t\in\mathbb R$, which implies $x_0<L^*$. Thus $w(t,x):=V(t,x)-\Phi^*(t,L^*-x)$ satisfies $w\geq 0$ in $\mathbb R \times(-\infty,L^*]$ and
\begin{equation*}\begin{cases}
  w_t-dw_{xx}-k^*(t)w_x-c(t,x)w=0 &\text{ for }(t,x)\in\mathbb R \times(-\infty,L^*),\\
  w(t,L^*)> 0 &\text{ for }t\in\mathbb R,
\end{cases}
\end{equation*}
with 
$$c(t,x):=\int_0^1\partial_uf(t,s\Phi^*+(1-s)V)ds\ \  \mbox{ a bounded function.}$$
  By the parabolic strong maximum principle, we deduce  $w(t,x)>0$ for $(t,x)\in\mathbb R \times(-\infty,L^*]$, which contradicts  $w(t_0,x_0)=0$. Hence, \eqref{eq-v>p} holds.

 \textbf{ Step 2.} We show that for any $y<L^*$, 
 \begin{equation}\label{eq-chi}
  \chi (y):=\inf_{(t,x)\in\mathbb{R}\times[y,L^*]}[V(t,x)-\Phi^*(t,L^*-x)]>0.
 \end{equation}
Clearly, $\chi(y)\geq 0$ for $y<L^*$. Assume to the contrary that \eqref{eq-chi} does not hold,  then there exists  $y_0\in(-\infty,L^*)$ such that 
\begin{equation*}
  \chi(y_0)=0.
\end{equation*}
Hence, we can find a sequence $\{(s_n,x_n)\}\subset \mathbb{R}\times [y_0,L^*]$ such that 
\begin{equation*}
  \lim_{n\to\infty}[V(s_n,x_n)-\Phi^*(s_n,L^*-x_n)]=0.
\end{equation*}
In view of Step 1, we have $|s_n|\to\infty$. Write $s_n=\tilde{k}_nT+\tilde{r}_n$ with $\tilde{k}_n\in\mathbb Z$, $\tilde r_n\in [0, T)$. Up to extraction of a subsequence if necessary, we may assume that $\tilde{r}_n\to\tilde{r}\in[0,T]$ and $x_n\to x_0\in[y_0,L^*]$. Set 
\begin{equation*}
  (V_n(t,x), \Gamma_n(t)):=(V(t+s_n-\tilde{r},x+x_n),\Gamma(t+s_n-\tilde{r})-x_n).
\end{equation*} 
 By applying the standard parabolic estimates and up to extraction of a subsequence,  we can argue similarly  to  Lemma \ref{lem-lim} to conclude that 
 \begin{equation*}
  (V_n(t,x), \Gamma_n(t))\to (V^\infty,\Gamma^\infty) \text{ in }C_{loc}^{(1+\alpha)/2,1+\alpha}(\overline{\Sigma}_\infty)\times C_{loc}^{1+\alpha/2}(\mathbb{R}),
 \end{equation*}
 where ${\Sigma}_\infty:=\{(t,x):t\in\mathbb{R},x<\Gamma^\infty(t)\}$ and $(V^\infty,\Gamma^\infty)$ satisfies
\begin{equation}\label{eq-Vinf}\begin{cases}
  V^\infty_t-dV^\infty_{xx}-k^*(t)V^\infty_x=f(t,V^\infty), V^\infty>\delta &\text{ for }(t,x)\in{\Sigma}_\infty,\\
  V^\infty(t,\Gamma^\infty(t))=\delta &\text{ for }t\in\mathbb R
\end{cases}
\end{equation}
and 
\begin{equation}\label{eq-Vin}
  V^\infty(\tilde{r},0)=\Phi^*(\tilde{r}, L^*-x_0).
\end{equation}
Moreover, $\Gamma_n(t)=\Gamma(t+s_n-\tilde{r})-x_n\geq \Gamma(r_*)-x_n $ leads to
 $\Gamma^\infty(t)+x_0\geq \Gamma(r_*)>L^*$ and \eqref{eq-v>p} yields
\begin{equation*}
  V^\infty(t,x)\geq \Phi^*(t,L^*-x_0-x)~~ \text{ for }(t,x)\in\mathbb{R} \times (-\infty, L^*-x_0].
\end{equation*}
In addition, it is easily seen that $\Phi^*(t,L^*-x_0-x)$ solves \eqref{eq-Vinf} with $\Gamma^\infty(t)$ replaced by $L^*-x_0$. Due to $L^\infty(t)>L^*-x_0>0$, we have
\begin{equation*}
  V^\infty(t,L^*-x_0)>\delta=\Phi^*(t,0) ~~\text{ for }t\in\mathbb{R}.
\end{equation*} 
Hence, we can apply the parabolic strong maximum principle to compare $V^\infty(t,x)$ with $\Phi^*(t,L^*-x_0-x)$ over $(-\infty,0]\times (-\infty,L^*-x_0)$ to conclude that 
\begin{equation*}
  V^\infty(t,x) >\Phi^*(t,L^*-x_0-x) \text{ for }t\leq \tilde{r},\ x<L^*-x_0,
\end{equation*}
which contradicts \eqref{eq-Vin}. Therefore, \eqref{eq-chi} holds for all $y< L^*$.

\textbf{Step 3.} We find a constant $L_0<L^*$ such that  $V(t,x)\geq \Phi^*(t,L^*-x+\epsilon)$ in $\mathbb{R}\times(-\infty,L_0]$ for all small $\epsilon>0$.

Since $f_u(\cdot,1)<0$ and $f$ is $C^1$ in $u$ uniformly for $t\in\mathbb R$,  there exists $\epsilon_0>0$ small enough such that  
\begin{equation}\label{eq-f_u}
  f_u(t,u)<0 ~~\text{ for }t\in\mathbb{R},\ u\in[1-\epsilon_0,1].
\end{equation} 
Owing to $\Phi^*(t,x) \to\infty$ as $x\to\infty$ uniformly in $t\in\mathbb R$,   we can find $L_0=L_0(\epsilon_0)<L^*$ with $L^*-L_0\gg 1$  such that 
\begin{equation}\label{eq-pl1}
  \Phi^*(t,L^*-x)\geq 1-\epsilon_0 ~~\text{ for }t\in\mathbb{R},\ x\leq L_0.
\end{equation}
By continuity,  there exists  $\epsilon\in(0,\epsilon_0)$ small enough (depending on $L^*-L_0$) such that
\begin{equation}\label{eq-pchi}
  \Phi^*(t,L^*-L_0+\epsilon)\leq \Phi^*(t,L^*-L_0)+\chi(L_0)~ \text{ for }t\in\mathbb R,
\end{equation}
where $\chi$ is defined by \eqref{eq-chi}.
Fix any such  $\epsilon$. For any $L<L_0$, consider the following auxiliary problem 
\begin{equation}\label{eq-wR}
  \begin{cases}
    w_t-dw_{xx}-k^*(t)w_x=f(t,w), &t>0, x\in(L,L_0),\\
    w(t,L)=1, ~w(t,L_0)=\Phi^*(t,L^*-L_0+\epsilon), &t>0,\\
    w(0,x)=1-\epsilon_0, &x\in[L,L_0].
  \end{cases}
\end{equation}
Clearly, $1$ and $1-\epsilon_0$ form a pair of upper and lower solutions of \eqref{eq-wR}.  It follows from the standard upper and lower solution arguments that \eqref{eq-wR} admits a maximal solution 
$$w_L\in C_{loc}^{1+\alpha/2,2+\alpha}((0,\infty)\times[L,L_0])\cap C_{loc}^{1+\alpha/2,2+\alpha}([0,\infty)\times(L,L_0))$$
 satisfying
\begin{equation}\label{eq-1pw}
 1-\epsilon_0\leq  w_L(t,x)\leq 1 ~~\text{ for }t>0,\ x\in[L,L_0].
\end{equation}
By applying the usual comparison principle, we know that $w_L$ increases in $L$ in the sense that $L_2<L_1$ implies $w_{L_2}(t,x)\leq w_{L_1}(t,x)$. Hence, the pointwise limit
\begin{equation*}
   w_\infty(t,x):=\lim_{L\to\infty}w_L(t,x) \ \mbox{ for } t>0,x\in(-\infty,L_0]
 \end{equation*}
 is well-defined and  $1-\epsilon_0\leq w_\infty\leq 1$. By applying  the standard parabolic estimates, we conclude that the above convergence also holds in $C_{loc}^{1,2}([0,\infty)\times(-\infty,L_0))$,  and thus $w_\infty$ solves \eqref{eq-wR} in $[0,\infty) \times(-\infty,L_0]$  except for the first equation in the second line.
Since $1-\epsilon_0$ is a  lower solution and $f$  is $T$-periodic in $t$,  $w_\infty(t+kT,x)$ is nondecreasing  in $k\in\mathbb R$. Hence, the pointwise limit
\begin{equation*}
  \tilde{w}(t,x)=\lim_{k\to\infty}w_\infty(t+kT,x)
\end{equation*}
is well-defined and $\tilde{w}(t,x)$ is $T$-periodic  in $t$ with $1-\epsilon_0\leq \tilde{w}\leq 1$.   Similar to before,  we deduce that  $\tilde{w}$ solves
\begin{equation}\label{eq-wtild}
  \begin{cases}
    \tilde{w}_t-d\tilde{w}_{xx}-k^*(t)\tilde{w}_x=f(t,\tilde{w}), ~1-\epsilon_0\leq \tilde{w}\leq 1, &t\in\mathbb R,\  x\in(-\infty,L_0),\\
    \tilde{w}(t,L_0)=\Phi^*(t,L^*-L_0+\epsilon), &t\in\mathbb R,\\
    \tilde{w}(t,x)=\tilde{w}(t+T,x), &t\in\mathbb R,\  x\in(-\infty,L_0].
  \end{cases}
\end{equation}
We claim that $\tilde{w}(t,-\infty)=1$. Indeed, for any sequence $\{x_n\}$  satisfying $x_n\to-\infty$ as $n\to\infty$, if we  denote 
\[\tilde{w}_n(t,x):=\tilde{w}(t,x+x_n) \text{ for } t\in\mathbb{R}, x\in(-\infty, L_0-x_n],\]
then $\tilde{w}_n(t,x)$ satisfies \eqref{eq-wtild} with $L_0$ replaced by $L_0-x_n$. Hence, by using the standard parabolic estimates and up to extraction of a subsequence, we know that there exists a function $\hat{w}$ such that 
\[\tilde{w}_n(t,x)\to\hat{w}(t,x) \text{  as } n\to\infty \text{ locally uniformly in } C^{1, 2}(\mathbb{R}\times\mathbb{R}).\]
Moreover, $\hat{w}(t,x)$ satisfies 
\[ \begin{cases}
  \hat{w}_t-d\hat{w}_{xx}-k^*(t)\hat{w}_x=f(t,\hat{w}), ~1-\epsilon_0\leq \hat{w}\leq 1, &t\in\mathbb R, x\in\mathbb{R},\\
  \hat{w}(t,x)=\hat{w}(t+T,x), &t\in\mathbb R, x\in\mathbb{R}.
\end{cases}
\]
The usual comparison principle yields that $\hat{w}(t,x)\geq p(t)$ with $p(t)$ solving
\[p_t=f(t,p),~t>0;~~~p(0)=1-\epsilon_0.\]
Due to $f(\cdot,u)>0$ for $u\in(0,1)$, we know that $p(t+kT)\to 1$  locally uniformly in $t\in\mathbb{R}$ as the integer $ k\to\infty$. Therefore, 
\[1\geq \hat{w}(t,x)=\lim_{k\to\infty}\hat{w}(t+kT,x)\geq \lim_{k\to\infty}p(t+kT)=1  \text{ locally uniformly in }(t,x)\in\mathbb{R}\times\mathbb R,\]
which implies that 
\[\tilde{w}(t,x+x_n)\to 1   \text{ locally uniformly in } C^{1, 2}(\mathbb{R}\times\mathbb{R}).\]
By the arbitrariness of  $\{x_n\}$, we conclude that $\tilde{w}(t,-\infty)=1$. The claim is thus proved.


We now set 
$$\Psi(t,x):=\Phi^*(t,L^*-x+\epsilon) ~\text{ for }t\in\mathbb R, x\leq L_0.$$
Clearly, $\Psi$ also satisfies \eqref{eq-wtild}. Thanks to \eqref{eq-pl1}, one has $\Psi(t,x)\geq 1-\epsilon_0$ for $x\leq L_0$. Therefore, we can compare $\Psi$ and $w_\infty$  to obtain
\begin{equation*}
  \Psi(t+kT,x)\geq  w_\infty(t+kT,x)~~\text{ for } t>-kT,\  x\in(-\infty,L_0],\  k\in\mathbb N.
\end{equation*}
Now, letting $k\to\infty$, one gets
\begin{equation}\label{eq-wp1}
  \Psi(t,x)\geq \tilde{w}(t,x)\geq 1-\epsilon_0~~\text{ for }t\in\mathbb{R},\ x\in(-\infty,L_0].
\end{equation}

We claim that 
\begin{equation*}
  \Psi(t,x)\equiv \tilde{w}(t,x)~~ \text{ for }t\in\mathbb{R},\  x\in(-\infty,L_0].
\end{equation*}
Denote
\begin{equation*}
  U(t,x):=\Psi(t,x)-\tilde{w}(t,x).
\end{equation*}
Then $U\geq 0$ satisfies 
\begin{equation}\label{eq-U}\begin{cases}
  U_t-dU_{xx}-k^*(t)U_x-f_u(t,\theta(t,x))U=0,&t\in\mathbb{R},\  x\in(-\infty,L_0),\\
  U(t,-\infty)=U(t,L_0)=0,&t\in\mathbb{R},\\
  U(t,x)=U(t+T,x), &t\in\mathbb{R},\ x\in(-\infty, L_0],
\end{cases}
\end{equation}
where $\theta(t,x)\in[\tilde{w}(t,x),\Psi(t,x)]\subset[1-\epsilon_0,1]$.
If the claim does not hold, then in view of \eqref{eq-wp1} and $U(t,-\infty)=0$, there exists $(t_0,x_0)\in\mathbb{R}\times(-\infty,L_0)$ such that 
$$U(t_0,x_0)=\max_{(t,x)\in\mathbb{R}\times(-\infty,L_0)}U(t,x)>0.$$
Since  $f_u(\cdot,u)<0$ for $u\in[1-\epsilon_0,1]$, using the first equation of \eqref{eq-U} we obtain
\begin{equation*}
  0=\big[U_t-dU_{xx}-k^*(t)U_x-f_u(t,\theta(t,x))U\big]\big|_{(t_0,x_0)}\geq -f_u(t,\theta(t_0,x_0))U(t_0,x_0)>0.
\end{equation*}
This contradiction proves the claim.

By \eqref{eq-v>p},\eqref{eq-chi} and \eqref{eq-pchi}, we have
\begin{eqnarray*}
  &&V(t,x)> \Phi^*(t,L^*-x)\geq 1-\epsilon_0 ~~\text{ for } (t,x)\in\mathbb R \times(-\infty,L_0],\\
  && V(t,L_0)\geq \Phi^*(t,L^*-L_0)+\chi (L_0)\geq  \Phi^*(t,L^*-L_0+\epsilon),
\end{eqnarray*}
 and thus we can compare $V$ with $w_\infty$  to deduce that 
\[ V(t,x)\geq w_\infty(t+kT,x)~\text{ for }t>-kT,\ x\in(-\infty,L_0],\  k\in\mathbb N.\]
Letting $k\to\infty$ we obtain 
\begin{equation}\label{eq-vw}
  V(t,x)\geq \tilde{w}(t,x)=\Phi^*(t,L^*-x+\epsilon) ~\text{ for }t\in\mathbb R,\ x\in(-\infty,L_0].
\end{equation}

\textbf{Step 4.} Completion of the proof.

By Step 2, we know that
\begin{equation*}
  \varrho:=\chi(L_0)>0.
\end{equation*}
Since $L^*<\Gamma(r_*)\leq \Gamma(t)$ for all $t\in\mathbb R$, we deduce from the  interior parabolic estimates that 
\begin{equation*}
  |V_x(t,x)|\leq C ~~\text{ for } t\in\mathbb R,~ L^*\leq x\leq (\Gamma(r_*)+L^*)/2,
\end{equation*}
and thus, for any $\epsilon_1\in(0,\varrho ]$ small enough, we have
\begin{eqnarray*}
 V(t,x)>V(t,L^*)-\varrho/2 ~~\text{ for }t\in\mathbb R,\  x\in[L^*,L^*+\epsilon_1].
\end{eqnarray*}
Similarly,  one has
\[\Phi^{*}(t,L^*-x+\epsilon_1)\leq \Phi^*(t,L^*-x)+\varrho/2 ~~\text{ for }t\in\mathbb R,\ x\in[L_0, L^*].\]
Recall that 
\begin{equation*}
V(t,x)-\Phi^*(t,L^*-x)\geq \chi (L_0)=\varrho  \text{ for }(t,x)\in\mathbb{R}\times[L_0,L^*].
\end{equation*}
As a consequence, for $t\in\mathbb{R}$ and $x\in[L_0,L^*]$, we readily have 
\begin{equation*}
  V(t,x)-\Phi^*(t,L^*-x+\epsilon_1)\geq \Phi^*(t,L^*-x)+\varrho - \Phi^*(t,L^*-x)-\varrho/2\geq \varrho/2>0.
\end{equation*}
Up to decreasing $\epsilon_1$ if necessary,  we may also assume
\begin{equation*}
  \Phi^*(t,L^*-x+\epsilon_1)<\Phi^*(t,0)+\varrho/2 ~~\text{ for } x\in[L^*,L^*+\epsilon_1],\ t\in\mathbb R.
\end{equation*}
Hence,  for $t\in\mathbb{R}$ and $x\in[L^*,L^*+\epsilon_1]$, we have
\begin{equation*}
  V(t,x)-\Phi^*(t,L^*-x+\epsilon_1)> V(t,L^*)-\varrho /2 - \Phi^*(t,0)-\varrho/2\geq0.
\end{equation*}
These inequalities, together  with \eqref{eq-vw}, finally give 
\begin{equation*}
  V(t,x)-\Phi^*(t,L^*-x+\epsilon_1)\geq0 \text{ for } t\in\mathbb{R},\ x\in(-\infty,L^*+\epsilon_1)
\end{equation*}
for all small $\epsilon_1\in(0,\epsilon)$, which violates the definition of $L^*$. Therefore, $L^*=\Gamma(r_*)$ and the lemma is thus  proved.
\end{proof}

\end{document}
